\section{Related work}\label{sec:related}

\subsection{LLM-evolutionary algorithm integration}

The automated heuristic design (AHD) literature has recently witnessed rapid progress, situated within the broader landscape of LLM-evolutionary algorithm (LLM-EA) integration~\cite{Wu2024LLMEA_Survey,Huang2024LLMEA_Roadmap}. Recent surveys characterize this emerging field along three axes: LLMs as operators (crossover, mutation), LLMs as fitness evaluators, and LLMs as search orchestrators~\cite{Chao2024LLMEvolutionary}. SolEvolve contributes primarily to the third category, extending orchestration beyond single-objective optimization to multi-modal mathematical discovery.

EvoPrompt~\cite{Chen2024EvoPrompt} and related works explore prompt optimization through evolutionary strategies, demonstrating that LLM outputs can be systematically improved via population-based search. The broader LLM agent literature~\cite{Wang2024LLMAgentSurvey,Xi2023LLMAgentRise} provides conceptual foundations for understanding SolEvolve's components: the Generator implements \textit{planning} and \textit{tool use}, the Verifier provides \textit{grounding}, and the Reflector performs \textit{self-reflection}---core capabilities identified as essential for autonomous agents.

\subsection{Automated heuristic design}

Within AHD specifically, SolSearch demonstrates that LLMs can iteratively edit SAT solver code, guided by curriculum prompts, to deliver double-digit improvements in PAR-2 scores~\cite{Zhang2025SolSearch}. FunSearch combines LLM reasoning with evolutionary search to discover new mathematical algorithms, including solutions to the cap set problem~\cite{RomeraParedes2024FunSearch}. Evolution of Heuristics (EoH-S) formalizes AHD as a supermodular set selection problem, using complementary portfolios to generalize across instance distributions~\cite{Liu2024EoHS}.

\subsection{Positioning of SolEvolve}

SolEvolve synthesizes insights from these paradigms while addressing their limitations through three key differentiators:

\begin{itemize}
\item \textbf{Multi-modal orchestration vs. single-domain optimization.} SolSearch focuses on single-solver optimization through curriculum-based code editing~\cite{Zhang2025SolSearch}. SolEvolve extends this to manage SAT solvers, ILP engines, and genetic algorithms within a unified framework, enabling cross-domain transfer of encoding strategies.

\item \textbf{Rigorous verification vs. fitness-based evaluation.} FunSearch targets mathematical function discovery through evolutionary code generation~\cite{RomeraParedes2024FunSearch}. SolEvolve adds rigorous verification through exhaustive validation, mathematical bound checking against authoritative databases (codetables.de), and formal UNSAT proofs for non-existence claims.

\item \textbf{Semantic feedback vs. metric-driven iteration.} Evolution of Heuristics (EoH) evolves both ``thoughts'' (natural language heuristic descriptions) and ``codes'' (executable implementations)~\cite{Liu2024EoH}. While EoH pioneered dual representation evolution, its feedback mechanism differs fundamentally from SolEvolve's. To clarify this distinction precisely: EoH's ``thoughts'' serve as \textit{prompts to the LLM for code generation}, guiding the mutation and crossover operators---the natural language descriptions help the LLM generate better code variants. However, the \textit{selection pressure} that determines which individuals survive is entirely \textit{scalar fitness-driven}: individuals with higher fitness scores (e.g., lower execution time, better solution quality) are selected regardless of \textit{why} they perform better. The ``thoughts'' do not inform selection; they only inform generation. This creates a disconnect: EoH cannot diagnose \textit{why} a promising heuristic failed or \textit{what structural property} caused stagnation---it only observes that fitness stopped improving.

In contrast, SolEvolve's Reflector provides \textit{diagnostic, causal-hypothesis-generating feedback} that directly informs both selection and generation. When the Evolver stagnates, the Reflector does not merely observe ``fitness plateau at 1,225.44 for 50 generations''; it analyzes the population structure to hypothesize \textit{why}: ``weight-12 codewords dominate the population, suggesting a local attractor in the weight-12 subspace.'' This diagnosis then generates targeted interventions: ``add symmetry-breaking constraints for columns 5--8 to escape the attractor'' or ``switch encoding from sequential counters to commander variables to alter the search landscape.'' This causal reasoning enables corrections that address root causes rather than symptoms, fundamentally distinguishing SolEvolve from fitness-driven approaches.
\end{itemize}
